% Научная публикация по проекту survivability
% Структура IMRAD; формулы — LaTeX; алгоритмы и интерфейсы — TikZ-макеты.
% Сборка: pdflatex doc/article.tex (два прохода)
\documentclass[12pt]{article}
\usepackage[T2A]{fontenc}
\usepackage[utf8]{inputenc}
\usepackage[russian]{babel}
\usepackage{amsmath, amssymb}
\usepackage[a4paper, margin=2.5cm]{geometry}
\usepackage{booktabs}
\setlength{\emergencystretch}{1.5em}
\usepackage{tikz}
\usetikzlibrary{arrows.meta, positioning, shapes.geometric, calc}
\DeclareMathOperator*{\argmax}{arg\,max}

\title{Методика и программный инструмент количественной оценки живучести\\
узла и системы компьютерной сети в условиях деструктивных воздействий}
\author{Программный комплекс survivability}
\date{}

\begin{document}
\maketitle

\begin{abstract}
Предложена методика и реализован открытый программный инструмент
количественной оценки живучести узла и системы компьютерной сети в условиях
деструктивных воздействий (случайные отказы, атаки, разрывы каналов связи).
Живучесть узла оценивается в марковской модели <<отказ--восстановление>> с
суммарной интенсивностью ущерба, образованной прореживанием пуассоновских
потоков воздействий; живучесть системы определяется как вероятность
достижения функцией эффективности допустимого порога по топологии сети и
вычисляется точным перебором состояний (до 20 узлов) либо методом
Монте-Карло с доверительным интервалом Уилсона; рациональный способ защиты
выбирается максимизацией живучести производной системы при ограничении на
бюджет, при этом сравнение вариантов выполняется на общем потоке случайных
чисел, что обеспечивает воспроизводимость и сопоставимость оценок.
Инструмент снабжён консольным интерфейсом, графической оболочкой
(редактор топологии, библиотека ролей узлов, экспорт схем SVG/DOT) и
набором модульных тестов. Приведены численные примеры: ранжирование
источников ущерба узла, сравнение топологий <<линия--треугольник>>,
сопоставление точного перебора и Монте-Карло, выбор рационального способа
защиты для линейной цепочки и для кольцевой сети из 22 сегментов.
\end{abstract}

\noindent\textbf{Ключевые слова:} живучесть; компьютерная сеть; марковская
модель; теория надёжности; метод Монте-Карло; доверительный интервал Уилсона;
критичность Бирнбаума; выбор способа защиты; цифровой двойник.

\medskip
\noindent\textbf{Keywords:} survivability; computer network; Markov model;
reliability theory; Monte Carlo method; Wilson score interval; Birnbaum
importance; protection selection; digital twin.

\section{Введение}

Живучесть --- свойство объекта сохранять способность выполнять заданные
функции при воздействии деструктивных факторов; в терминологии ГОСТ 27.002
она относится к составляющим надёжности наряду с безотказностью и
долговечностью~\cite{gost}. Для компьютерных сетей деструктивные факторы ---
это собственные случайные отказы оборудования, атаки (DDoS, внедрение
вредоносного кода, фишинг), радиоэлектронное подавление и разрывы каналов
связи. Количественная оценка живучести нужна при проектировании сегментов
сети, обосновании SLA, приоритизации мер защиты и проверке нормативов
устойчивости функционирования.

\begin{sloppypar}
Классический аппарат теории надёжности --- марковские модели
<<отказ--восстановление>>, пуассоновские потоки событий и показатели
структурной важности --- хорошо разработан для технических
систем~\cite{gnedenko,barlow,ventcel} и для приложений в информатике
\cite{trivedi,rausand}. Однако на практике оценки живучести сетей часто
получают раздельным учётом отказов и атак либо на единичных сценариях
<<ручного>> анализа; первый подход завышает живучесть из-за двойного
учёта, второй невоспроизводим и не даёт доверительных границ. Оценка
вероятностных показателей методом Монте-Карло воспроизводима, но сравнение
вариантов защиты по независимым прогонам смешивает различия вариантов с
случайностью потоков чисел; в статистике симуляций эту проблему решают
техникой общих (сопряжённых) потоков случайных чисел
\cite{glasserman,ross}.
\end{sloppypar}

Цель работы --- единый программный инструмент, в котором три
последовательные задачи формируют непрерывный поток вычислений:
(1)~оценка живучести узла $Q_i$ в марковской модели с совместным учётом
отказов и воздействий; (2)~оценка живучести системы $Q(S)$ по топологии
сети с показателями важности узлов; (3)~выбор рационального способа
защиты $g^{*}$ при ограничении на бюджет. Новизна решения: суммарная
интенсивность ущерба как результат прореживания пуассоновских потоков
(совместная динамика отказов и атак без двойного учёта); сопоставление
вариантов защиты на общем потоке случайных чисел с фиксированным зерном,
делающее выбор детерминированным и воспроизводимым; открытая реализация на
Python со стандартной библиотекой в ядре, графической оболочкой и
воспроизводимыми прогонами <<цифрового двойника>> сети.

Статья организована по схеме IMRAD: раздел~2 излагает методы (модель,
формулы, алгоритмы, архитектура), раздел~3 приводит численные результаты
на реальных входных данных инструмента, раздел~4 обсуждает ограничения
и формулирует выводы.

\section{Методы}

\subsection{Постановка задач}

Система описывается структурной моделью $S = \langle E, L, F \rangle$, где
$E = \{e_1, \dots, e_n\}$~--- множество узлов, $L$~--- множество
неориентированных связей, $F$~--- целевая функция. На узел $e_i$ действуют
деструктивные воздействия $D = \{d_1, \dots, d_m\}$ с интенсивностями
$\lambda_j$; каждое воздействие поражает узел с вероятностью
$p_{ij} = P(x_i = 0 \mid d_j)$, где $x_i \in \{0, 1\}$~--- состояние узла
(1~--- работоспособен, 0~--- отказ). За период оценивания $T$ требуется
найти:

\textbf{Задача 1} --- живучесть узла: $Q_i = P(x_i(T) = 1)$ при условии
начальной работоспособности $x_i(0) = 1$.

\textbf{Задача 2} --- живучесть системы: $Q(S) = P(W(X) \geq
W_{\text{доп}})$, где $X = (x_1, \dots, x_n) \in \Omega = \{0,1\}^n$~---
вектор состояний, $W(X) \in [0, 1]$~--- функция эффективности,
$W_{\text{доп}}$~--- допустимый порог деградации; а также показатели
важности узлов.

\textbf{Задача 3} --- рациональный способ защиты: из множества вариантов
$G$ со стоимостями $C(g)$ выбрать допустимый по бюджету $C_{\text{доп}}$
вариант с максимальной живучестью производной системы $S_g$.

\subsection{Модель деструктивных воздействий и живучесть узла}

Собственные отказы узла и каждое воздействие образуют пуассоновские потоки
событий~\cite{ventcel}. Реализация воздействия $d_j$ прориживает поток
с вероятностью $p_{ij}$, давая пуассоновский поток поражений с
интенсивностью $\lambda_j p_{ij}$; сумма независимых пуассоновских потоков
--- пуассоновский поток, поэтому общая интенсивность ущерба узла
\begin{equation}
\Lambda_i = \lambda_i^{(0)} + \sum_{j=1}^{m} \lambda_j\, p_{ij},
\label{eq:lambda}
\end{equation}
где $\lambda_i^{(0)}$~--- интенсивность собственных отказов. Совместная
динамика отказов и атак учитывается одной интенсивностью, что не завышает
оценку в отличие от раздельного учёта с последующим <<перемножением>>
выживаемости.

Функционирование узла описывается однородной марковской цепью с
непрерывным временем и двумя состояниями (рис.~\ref{fig:markov}):
переход $1 \to 0$~--- с интенсивностью $\Lambda_i$, обратный переход
(восстановление)~--- с интенсивностью $\mu \geq 0$ (экспоненциальное время
восстановления; $\mu = 0$ соответствует его отсутствию). Генератор цепи:
$A = \bigl(\begin{smallmatrix} -\Lambda_i & \Lambda_i \\ \mu & -\mu
\end{smallmatrix}\bigr)$.

\begin{figure}[htb]
\centering
\begin{tikzpicture}[state/.style={circle, draw, minimum size=1.2cm,
                                  font=\large}, node distance=4.2cm]
\node[state] (one) {1};
\node[state, right=of one] (zero) {0};
\node[below=1mm of one] {\footnotesize работоспособен};
\node[below=1mm of zero] {\footnotesize отказ};
\draw[-{Stealth}, bend left=22] (one) to node[above] {$\Lambda_i$} (zero);
\draw[-{Stealth}, bend left=22] (zero) to node[below] {$\mu$} (one);
\end{tikzpicture}
\caption{Марковская модель <<отказ--восстановление>> узла}
\label{fig:markov}
\end{figure}

Решение уравнения Колмогорова с начальным условием $P_1(0) = 1$ для
случая без восстановления ($\mu = 0$) даёт вероятность безотказной работы
\begin{equation}
Q_i(T) = e^{-\Lambda_i T},
\label{eq:cold}
\end{equation}
а при $\mu > 0$~--- живучесть узла
\begin{equation}
Q_i(T) = \frac{\mu}{\Lambda_i + \mu}
+ \frac{\Lambda_i}{\Lambda_i + \mu}\, e^{-(\Lambda_i + \mu) T}.
\label{eq:main}
\end{equation}
Первое слагаемое~--- стационарный коэффициент готовности, второе~---
вероятность того, что за время $T$ узел ни разу не сменил состояние:
\begin{equation}
K = \lim_{T \to \infty} Q_i(T) = \frac{\mu}{\Lambda_i + \mu}.
\label{eq:steady}
\end{equation}
Частные случаи согласованы: $T \to 0 \Rightarrow Q_i \to 1$;
$\mu = 0$ сводится к~\eqref{eq:cold}; $\Lambda_i = 0$ даёт $Q_i = 1$.
Формулы~\eqref{eq:cold}--\eqref{eq:steady}~--- стандартные результаты
теории надёжности двухсостоятельной системы с
восстановлением~\cite{gnedenko,barlow,trivedi}.

Для ранжирования угроз каждому источнику ущерба сопоставляется его доля
в общей интенсивности:
\begin{equation}
s_j = \frac{\lambda_j p_{ij}}{\Lambda_i}, \qquad
s_0 = \frac{\lambda_i^{(0)}}{\Lambda_i}, \qquad
s_0 + \sum_{j=1}^{m} s_j = 1
\quad (\text{при } \Lambda_i > 0).
\label{eq:shares}
\end{equation}

\subsection{Живучесть системы}

Состояния узлов принимаются независимыми с $P(x_i = 1) = Q_i$
по~\eqref{eq:main}, поэтому вероятность вектора состояний факторизуется:
\begin{equation}
P(X) = \prod_{i=1}^{n} Q_i^{x_i}\, (1 - Q_i)^{1 - x_i}.
\label{eq:px}
\end{equation}
Топология $L$ входит в оценку через функцию эффективности. Реализованы
два вида целевой функции $F$:

--- доля работоспособных узлов (связи не учитываются):
\begin{equation}
W(X) = \frac{1}{n} \sum_{i=1}^{n} x_i;
\label{eq:wfraction}
\end{equation}

--- доля пар узлов, между которыми существует путь по работоспособным
узлам и связям (по компонентам связности $s_k$ работоспособной части
графа):
\begin{equation}
W(X) = \frac{\sum_k \dbinom{s_k}{2}}{\dbinom{n}{2}},
\label{eq:wconn}
\end{equation}
при $n = 1$ полагается $W = x_1$. Для~\eqref{eq:wconn} компоненты связности
выделяются системой непересекающихся множеств за $O(|L|)$ на вектор
состояний.

Живучесть системы~--- вероятность допустимого состояния:
\begin{equation}
Q(S) = P\!\left(W(X) \geq W_{\text{доп}}\right)
     = \sum_{X:\, W(X) \geq W_{\text{доп}}} P(X).
\label{eq:qs}
\end{equation}

При $n \leq 20$ сумма~\eqref{eq:qs} вычисляется прямым перебором всех
$2^n$ векторов (вырожденные случаи $Q_i \in \{0, 1\}$ обрабатываются
естественно). При $n > 20$ выполняется $N$ независимых прогонов, в каждом
состояние узла генерируется по схеме Бернулли с вероятностью $Q_i$;
оценка $\hat Q(S) = N^{*}/N$, где $N^{*}$~--- число прогонов с допустимым
состоянием. Точность задаётся доверительным интервалом Уилсона с
доверительной вероятностью 95\,\% ($z = 1{,}96$), корректным и при малых
$N^{*}$, в отличие от интервала на основе нормальной
аппроксимации~\cite{wilson,cramer}:
\begin{equation}
\hat Q(S) \in
\left[
\frac{\hat p + \frac{z^2}{2N} - z\sqrt{\frac{\hat p(1-\hat p)}{N}
+ \frac{z^2}{4N^2}}}{1 + \frac{z^2}{N}}
\;;\;
\frac{\hat p + \frac{z^2}{2N} + z\sqrt{\frac{\hat p(1-\hat p)}{N}
+ \frac{z^2}{4N^2}}}{1 + \frac{z^2}{N}}
\right], \qquad \hat p = \frac{N^{*}}{N}.
\label{eq:wilson}
\end{equation}
Статистическая теория оценок по прогонам изложена в~\cite{cramer,metropolis}.

Для приоритизации узлов под защиту используется критичность
Бирнбаума~\cite{birnbaum}~--- разность условных живучестей:
\begin{equation}
I_i = Q(S \mid x_i = 1) - Q(S \mid x_i = 0),
\label{eq:birnbaum}
\end{equation}
\begin{equation}
Q(S \mid x_i = 1) = \frac{1}{Q_i}\!\!
\sum_{\substack{X:\, x_i = 1 \\ W(X) \geq W_{\text{доп}}}}\!\! P(X),
\qquad
Q(S \mid x_i = 0) = \frac{1}{1 - Q_i}\!\!
\sum_{\substack{X:\, x_i = 0 \\ W(X) \geq W_{\text{доп}}}}\!\! P(X).
\label{eq:conditional}
\end{equation}
Дополнительно вычисляется потенциал повышения живучести
$\Delta_i^{+} = Q(S \mid x_i = 1) - Q(S)$~--- прирост $Q(S)$ от идеального
узла. При точном переборе суммы~\eqref{eq:conditional} накапливаются за один
проход по $\Omega$; при Монте-Карло~--- $2\,000$ условных прогонов на узел
(узел $i$ фиксируется работоспособным или отказавшим, остальные
генерируются случайно).

\subsection{Выбор рационального способа защиты}

Каждый вариант защиты $g \in G$ описывается изменёнными параметрами узлов
(сниженные $\lambda_i^{(0)}$, повышенное $\mu$, пониженные $p_{ij}$)
и/или добавленными резервными связями, а также стоимостью $C(g) \geq 0$.
Вариант порождает производную систему $S_g$, для которой $Q_g(S)$
оценивается формулами~\eqref{eq:px}--\eqref{eq:wilson} тем же способом.
Рациональный способ защиты:
\begin{equation}
g^{*} = \argmax_{g \in G:\; C(g) \leq C_{\text{доп}}} Q_g(S);
\label{eq:argmax}
\end{equation}
если допустимых вариантов нет, решение отсутствует; при равенстве
$Q_g(S)$ выбирается вариант с меньшей стоимостью, затем с меньшим именем
(детерминированное правило разрыва). Постановка согласуется с подходами
оптимизации надёжности при ограничениях на
затраты~\cite{rausand,barlow}.

Ключевой элемент сопоставимости~--- \emph{общий поток случайных чисел}:
все оценки (базовая система и все варианты) выполняются последовательно
над одним генератором с фиксированным зерном $seed$ (техника сопряжённых
потоков~\cite{glasserman,ross}). Оценки $\hat Q_g(S)$ зависимы, их
различия отражают различия вариантов, а не случайность потоков; при
фиксированном зерне результат сравнения детерминирован и любой проверяющий
воспроизводит его одной командой. Обоснование решения сопровождается
производными показателями:
\begin{equation}
\Delta Q = Q_{g^{*}}(S) - Q_{\text{base}}(S), \qquad
\frac{\Delta Q}{Q_{\text{base}}}, \qquad
\frac{\Delta Q}{C(g^{*})},
\label{eq:deltaq}
\end{equation}
где $Q_{\text{base}}(S)$~--- живучесть системы без защиты, а
$\Delta Q / C(g^{*})$~--- эффективность затрат (прирост живучести на
единицу стоимости).

\subsection{Архитектура программного инструмента}

Инструмент реализован на Python~3.11+; ядро расчётов использует только
стандартную библиотеку, что упрощает аудит и внедрение. Архитектура
(рис.~\ref{fig:arch}) состоит из шести слоёв: модель данных и валидация
(\texttt{model.py}); расчёт живучести узла (\texttt{calc.py}); оценка
живучести системы, Монте-Карло и важность (\texttt{system\_calc.py});
выбор защиты с общим потоком случайных чисел (\texttt{protection.py});
кодек проекта и экспорт топологии в SVG/DOT без зависимости от
графической библиотеки (\texttt{project\_codec.py}); консольный
(\texttt{cli.py}) и графический (\texttt{gui.py}, PyQt6) интерфейсы поверх
общего ядра. Файл проекта~--- JSON, совместимый между CLI и GUI: метаданные
оболочки хранятся в служебной секции \texttt{gui}, которую консольная
утилита игнорирует. Набор включает 73+ модульных теста (unittest).

\begin{figure}[htb]
\centering
\begin{tikzpicture}[box/.style={draw, rounded corners, align=center,
                               minimum height=1.05cm, text width=3.4cm,
                               font=\small}, node distance=0.35cm]
\node[box] (model) {\texttt{model.py}\\модель и валидация};
\node[box, right=of model] (calc) {\texttt{calc.py}\\$Q_i$, $K$, вклады $s_j$};
\node[box, right=of calc] (syscalc) {\texttt{system\_calc.py}\\$Q(S)$: перебор / Монте-Карло, $I_i$};
\node[box, right=of syscalc] (prot) {\texttt{protection.py}\\варианты, бюджет, $g^{*}$};
\node[box, below=0.6cm of model] (codec) {\texttt{project\_codec.py}\\проект $\leftrightarrow$ JSON, SVG/DOT};
\node[box, below=0.6cm of calc] (io) {\texttt{io.py}\\ввод/вывод, отчёты};
\node[box, below=0.6cm of syscalc] (ui) {интерфейсы:\\ \texttt{cli.py} (\texttt{\_\_main\_\_.py});\\ \texttt{gui.py} (PyQt6)};
\draw[-{Stealth}, gray] (calc) -- (model);
\draw[-{Stealth}, gray] (syscalc) -- (calc);
\draw[-{Stealth}, gray] (prot) -- (syscalc);
\draw[-{Stealth}, gray] (codec) -- (model);
\draw[-{Stealth}, gray] (io) -- (codec);
\draw[-{Stealth}, gray] (ui) -- (io);
\end{tikzpicture}
\caption{Архитектура программного комплекса; стрелки~--- направление
зависимостей модулей}
\label{fig:arch}
\end{figure}

Графическая оболочка (макет~--- рис.~\ref{fig:gui-editor}) содержит
редактор узлов с полями $\lambda_i^{(0)}$, $\mu$, $T$ и таблицей
воздействий ($d_j$, $\lambda_j$, $p_{ij}$), канвас топологии с
перетаскиванием узлов (координаты сохраняются в файле проекта), палитру
ролей узлов (сервер, маршрутизатор, АРМ, СХД; параметры типового
устройства подставляются автоматически), редактор вариантов защиты и
панель результатов (макет~--- рис.~\ref{fig:gui-results}). Вся валидация
выполняется через конструкторы моделей: при ошибке расчёт не запускается,
а текст ошибки показывается диалогом. Результаты экспортируются в JSON/CSV
теми же функциями, что использует CLI, а топология~--- в SVG с раскраской
узлов по $Q_i$ или важности $I_i$.

\begin{figure}[htb]
\centering
\begin{tikzpicture}[font=\scriptsize]
\draw[rounded corners, thick] (0,0) rectangle (12.4,7.0);
\fill[gray!25] (0,6.7) rectangle (12.4,7.0);
\node[anchor=west] at (0.15,6.85) {survivability --- редактор проекта};
% палитра ролей
\draw (0.15,0.55) rectangle (2.35,6.55);
\node[anchor=north] at (1.25,6.55) {\textbf{Роли}};
\foreach \y/\t in {5.95/{Сервер}, 5.35/{Маршрутизатор}, 4.75/{АРМ}, 4.15/{СХД}}{
  \node[draw, rounded corners, minimum width=1.95cm, minimum height=0.4cm]
    at (1.25,\y) {\t};}
% канвас
\draw[fill=white] (2.55,0.55) rectangle (9.15,6.55);
\node[anchor=north west] at (2.65,6.5) {\textbf{Канвас топологии}};
\node[draw, rounded corners, minimum width=1.75cm, minimum height=0.62cm]
  (coi) at (4.1,4.7) {ЦОИ};
\node[draw, rounded corners, minimum width=1.75cm, minimum height=0.62cm]
  (db) at (6.6,4.7) {Сервер БД};
\node[draw, rounded corners, minimum width=1.75cm, minimum height=0.62cm]
  (arm) at (5.35,2.7) {АРМ оператора};
\draw[-] (coi) -- (db);
\draw[-] (db) -- (arm);
% панель свойств
\draw (9.35,0.55) rectangle (12.25,6.55);
\node[anchor=north, align=center] at (10.8,6.55) {\textbf{Свойства узла}};
\foreach \y/\t in {5.9/{$\lambda_i^{(0)} = 0{,}01$},
                   5.35/{$\mu = 4{,}0$},
                   4.8/{$T = 1{,}0$},
                   4.25/{воздействий: 3}}{
  \node[anchor=west] at (9.45,\y) {\t};}
% нижняя панель
\draw (0.15,0.12) rectangle (12.25,0.5);
\node at (6.2,0.31) {Рассчитать $Q(S)$ \; $\vert$ \; Сравнить варианты
защиты \; $\vert$ \; Экспорт SVG/DOT};
\end{tikzpicture}
\caption{Макет основного окна графической оболочки: палитра ролей,
канвас топологии, панель свойств узла}
\label{fig:gui-editor}
\end{figure}

\begin{figure}[htb]
\centering
\fbox{\begin{tabular}{@{}lccc@{}}
\multicolumn{4}{l}{\textbf{Сравнение способов защиты},
$C_{\text{доп}} = 50$,\; $Q_{\text{base}}(S) = 0{,}701266$}\\[2pt]
\toprule
Вариант & $C(g)$ & Допустим & $Q_g(S)$ \\
\midrule
Экранирование ЦОИ & 25 & да & 0{,}712178 \\
\textbf{Резервный канал ЦОИ--АРМ} & 40 & да & \textbf{0{,}897544} \\
Комплексная защита & 70 & нет & 0{,}949679 \\
\bottomrule
\end{tabular}}
\caption{Макет панели результатов выбора способа защиты (пример
разд.~\ref{sec:protection-line})}
\label{fig:gui-results}
\end{figure}

\subsection{Алгоритм расчёта}

Блок-схема потока вычислений задач 2--3 приведена на
рис.~\ref{fig:flow}. Точный перебор и Монте-Карло различаются только
способом оценки суммы~\eqref{eq:qs}; показатели важности вычисляются
за один дополнительный проход (перебор) или условными прогонами
(Монте-Карло).

\begin{figure}[htb]
\centering
\begin{tikzpicture}[blk/.style={draw, rounded corners, align=center,
                               minimum height=0.8cm, inner sep=4pt,
                               text width=4.4cm, font=\small},
                    dec/.style={draw, diamond, aspect=2.6, align=center,
                                inner sep=1pt, font=\small},
                    lbl/.style={font=\scriptsize},
                    node distance=0.5cm]
\node[blk] (in) {Вход: \texttt{SystemInput} --- узлы $E$, связи $L$,
$W_{\text{доп}}$, \texttt{target}, $N$, \texttt{seed}};
\node[blk, below=of in] (qi) {Живучести узлов $Q_i$
(формула~\eqref{eq:main}): $\Lambda_i$, $\mu = 0$ или $\mu > 0$};
\node[blk, below=of qi] (adj) {Списки смежности; выбор $W(X)$:
\texttt{node\_fraction} \eqref{eq:wfraction} или
\texttt{connectivity} \eqref{eq:wconn}};
\node[dec, below=of adj] (d) {$n \leq 20$?};
\node[blk, left=1.2cm of d, text width=3.6cm] (ex) {Точный перебор: маски
$X = 0..2^n - 1$, сумма $P(X)$ по допустимым};
\node[blk, right=1.2cm of d, text width=3.6cm] (mc) {Монте-Карло: $N$
прогонов Бернулли$(Q_i)$; $\hat Q = N^{*}/N$; ДИ Уилсона
\eqref{eq:wilson}};
\node[blk, below=2.6cm of d, text width=6.2cm] (imp) {Важность узлов $I_i$
\eqref{eq:birnbaum}: перебор~--- накопление сумм за один проход;
Монте-Карло~--- $2\,000$ условных прогонов на узел};
\node[blk, below=of imp] (out) {Выход: оценка $Q(S)$ (способ, ДИ),
список \texttt{NodeImportance}};
\draw[-{Stealth}] (in) -- (qi);
\draw[-{Stealth}] (qi) -- (adj);
\draw[-{Stealth}] (adj) -- (d);
\draw[-{Stealth}] (d) -- node[lbl, above] {да} (ex);
\draw[-{Stealth}] (d) -- node[lbl, above] {нет} (mc);
\draw[-{Stealth}] (ex.south) |- ($(imp.west) + (0,0.25)$) -- (imp.west);
\draw[-{Stealth}] (mc.south) |- ($(imp.east) + (0,0.25)$) -- (imp.east);
\draw[-{Stealth}] (imp) -- (out);
\end{tikzpicture}
\caption{Блок-схема оценки живучести системы и важности узлов}
\label{fig:flow}
\end{figure}

\section{Результаты}

Все численные результаты получены релизным интерфейсом инструмента
(\texttt{python -m survivability {-}-file <проект>.json}, Python 3.14,
64-разрядная Windows, пользовательский класс ПК) на входных данных из
поставки инструмента и воспроизводимы; для сценариев Монте-Карло зерно
\texttt{seed} зафиксировано во входных файлах.

\subsection{Живучесть узла и ранжирование источников ущерба}

Рассмотрим сервер маршрутизации ($T = 24$~ч, $\lambda_i^{(0)} = 0{,}02$,
$\mu = 2$) с тремя классами воздействий (табл.~\ref{tab:node}). Общая
интенсивность ущерба $\Lambda_i = 0{,}845$; поскольку $(\Lambda_i + \mu)T
\approx 68{,}3$, транзиентный член формулы~\eqref{eq:main} пренебрежимо мал
и $Q_i(T) \approx K = 2 / 2{,}845 = 0{,}702988$: живучесть сервера за
сутки определяется стационарным коэффициентом готовности, а период $T$
уже не играет роли.

\begin{table}[htb]
\centering
\caption{Вклады источников ущерба узла (сервер маршрутизации)}
\label{tab:node}
\begin{tabular}{@{}lcclr@{}}
\toprule
Источник & $\lambda_j$ & $p_{ij}$ & $\lambda_j p_{ij}$ & Доля $s_j$ \\
\midrule
Собственные отказы & 0{,}02 & --- & 0{,}020 & 2{,}37\,\% \\
DDoS-атака & 1{,}5 & 0{,}4 & 0{,}600 & 71{,}01\,\% \\
Внедрение вредоносного кода & 0{,}05 & 0{,}9 & 0{,}045 & 5{,}33\,\% \\
Отказ канала связи & 0{,}3 & 0{,}6 & 0{,}180 & 21{,}30\,\% \\
\midrule
Итого ($\Lambda_i$) & \multicolumn{2}{c}{} & 0{,}845 & 100\,\% \\
\bottomrule
\end{tabular}
\end{table}

Разбивка~\eqref{eq:shares} показывает, что $71\,\%$ ущерба даёт DDoS-атака
и лишь $2{,}4\,\%$~--- собственные отказы: меры защиты, нацеленные на
отказоустойчивое оборудование, почти не меняют $Q_i$, эффективна защита от
DDoS. Таким образом, модель переводит спор о <<главной угрозе>> в область
вычислимых долей $s_j$.

\subsection{Живучесть системы: ручной расчёт и сверка с программой}

Возьмём систему из трёх узлов в линию <<ЦОИ --- Сервер БД --- АРМ
оператора>> ($T = 1$, \texttt{target = connectivity}, $W_{\text{доп}} =
1/3$). Живучести узлов по формуле~\eqref{eq:main} (табл.~\ref{tab:nodes}):
например, для ЦОИ
\[
Q_1 = \frac{4}{4{,}78} + \frac{0{,}78}{4{,}78}\, e^{-4{,}78}
    = 0{,}838190.
\]

\begin{table}[htb]
\centering
\caption{Исходные данные и живучести узлов линейной цепочки}
\label{tab:nodes}
\begin{tabular}{@{}lcrrcr@{}}
\toprule
Узел & $\lambda_i^{(0)}$ & $\sum_j \lambda_j p_{ij}$ & $\Lambda_i$ &
$\mu$ & $Q_i$ \\
\midrule
ЦОИ & 0{,}01 & 0{,}770 & 0{,}780 & 4{,}0 & 0{,}838190 \\
Сервер БД & 0{,}02 & 0{,}825 & 0{,}845 & 2{,}0 & 0{,}720254 \\
АРМ оператора & 0{,}05 & 0{,}240 & 0{,}290 & 1{,}0 & 0{,}837076 \\
\bottomrule
\end{tabular}
\end{table}

Полный перебор $2^3 = 8$ состояний (табл.~\ref{tab:states}) даёт
\[
Q(S) = 0{,}098359 + 0{,}097556 + 0{,}505351 = 0{,}701266.
\]
Программа на том же входе возвращает $Q(S) = 0{,}701266$; расхождения с
ручным расчётом отсутствуют до шестого знака. Ключевая строка
табл.~\ref{tab:states}~--- состояние $(101)$: оба крайних узла
работоспособны, но путь между ними проходит через отказавший Сервер БД,
поэтому $W = 0$ и состояние недопустимо. Это и есть вклад топологии $L$
в формулу~\eqref{eq:qs}: мера $P(X)$ зависит только от узлов, а
допустимость состояния~--- только от геометрии сети.

\begin{table}[htb]
\centering
\caption{Перебор состояний линейной цепочки ($W_{\text{доп}} = 1/3$)}
\label{tab:states}
\begin{tabular}{@{}clrrc@{}}
\toprule
$X$ & Работоспособны & $P(X)$ & $W(X)$ & Допустим \\
\midrule
(000) & --- & 0{,}007375 & 0{,}000000 & нет \\
(100) & ЦОИ & 0{,}038202 & 0{,}000000 & нет \\
(010) & Сервер БД & 0{,}018988 & 0{,}000000 & нет \\
(110) & ЦОИ + Сервер БД & 0{,}098359 & 0{,}333333 & да \\
(001) & АРМ оператора & 0{,}037891 & 0{,}000000 & нет \\
(101) & ЦОИ + АРМ оператора & 0{,}196278 & 0{,}000000 & нет \\
(011) & Сервер БД + АРМ оператора & 0{,}097556 & 0{,}333333 & да \\
(111) & все & 0{,}505351 & 1{,}000000 & да \\
\bottomrule
\end{tabular}
\end{table}

Критичности Бирнбаума (табл.~\ref{tab:importance}) показывают, что Сервер
БД~--- точка сочленения линии: при его гарантированном отказе $Q(S \mid
x_2 = 0) = 0$ (сеть распадается на два изолированных узла) и $I_2 =
0{,}973637 \approx 1$. Крайние узлы взаимозаменяемы, их критичности почти
совпадают ($0{,}117346$ и $0{,}116544$).

\begin{table}[htb]
\centering
\caption{Важность узлов линейной цепочки}
\label{tab:importance}
\begin{tabular}{@{}lrrr@{}}
\toprule
Узел & $Q(S \mid x_i = 1)$ & $Q(S \mid x_i = 0)$ & $I_i$ \\
\midrule
ЦОИ & 0{,}720254 & 0{,}602908 & 0{,}117346 \\
Сервер БД & 0{,}973637 & 0{,}000000 & \textbf{0{,}973637} \\
АРМ оператора & 0{,}720254 & 0{,}603710 & 0{,}116544 \\
\bottomrule
\end{tabular}
\end{table}

\subsection{Влияние топологии и целевой функции}

Те же три узла ($Q_i$ не менялись) с разной структурой связей дают разную
живучесть (табл.~\ref{tab:topologies}). Добавление одной резервной связи
<<ЦОИ--АРМ>> (треугольник) делает допустимым состояние $(101)$ и поднимает
$Q(S)$ с $0{,}701266$ до $0{,}897544$ ($+28\,\%$); при этом критичности
выравниваются до $0{,}27$--$0{,}35$, поскольку отказ любого одиночного
узла оставляет остальные две пары связными. Смена целевой функции на долю
работоспособных узлов без связей даёт $Q(S) = 0{,}992625$~--- почти
вероятность <<хотя бы один узел жив>>: одинокий работоспособный узел уже
допустим ($W = 1/3$), и вклад топологии обнуляется. Для второго примера
(два резервных сервера за клиентом, \texttt{connectivity}) программа даёт
$Q(S) = 0{,}752894$ с доминирующей критичностью клиента $I = 0{,}899433$.

\begin{table}[htb]
\centering
\caption{Живучесть и важность узлов при разной структуре связей
(одни и те же $Q_i$)}
\label{tab:topologies}
{\small
\begin{tabular}{@{}lcrrrr@{}}
\toprule
Структура & \texttt{target} & $Q(S)$ & $I$(ЦОИ) & $I$(Сервер БД) &
$I$(АРМ) \\
\midrule
Линия (2 связи) & connectivity & 0{,}701266 & 0{,}117346 &
\textbf{0{,}973637} & 0{,}116544 \\
Треугольник (3 связи) & connectivity & 0{,}897544 & 0{,}351515 &
0{,}272008 & 0{,}351024 \\
Без связей & node\_fraction & 0{,}992625 & 0{,}045577 & 0{,}026363 &
0{,}045266 \\
\bottomrule
\end{tabular}}
\end{table}

Три наблюдения обобщают пример: (1)~топология~--- самостоятельный
множитель живучести, сравнимый по силе с надёжностью узлов; (2)~важность
узла~--- свойство его положения в сети, а не только собственной
надёжности $Q_i$; (3)~выбор \texttt{target} должен отвечать на вопрос,
что для системы означает работоспособность.

\subsection{Точный перебор и метод Монте-Карло}

Для линейной цепочки ($Q(S) = 0{,}701266$) выполним Монте-Карло с
\texttt{seed} $= 42$ при разном числе прогонов $N$ (табл.~\ref{tab:mc}).
Оценки $\hat Q(S)$ сходятся к точному значению, а интервал Уилсона при
$N = 100\,000$ ($\pm 0{,}003$) уже достаточен для инженерных выводов.

\begin{table}[htb]
\centering
\caption{Монте-Карло против точного перебора (линейная цепочка,
$Q(S) = 0{,}701266$)}
\label{tab:mc}
\begin{tabular}{@{}rrlr@{}}
\toprule
$N$ & $\hat Q(S)$ & ДИ 95\,\% & $|\hat Q - Q|$ \\
\midrule
1\,000 & 0{,}701000 & [0{,}671899; 0{,}728562] & 0{,}000266 \\
10\,000 & 0{,}699600 & [0{,}690540; 0{,}708507] & 0{,}001666 \\
100\,000 & 0{,}702690 & [0{,}699849; 0{,}705515] & 0{,}001424 \\
\bottomrule
\end{tabular}
\end{table}

Время точного перебора растёт как $2^n$ (табл.~\ref{tab:scaling}): при
$n = 20$ перебор занимает $11{,}4$~с, что допустимо для интерактивного
расчёта, но уже при $n = 21$ время удваивается. Монте-Карло при $n = 20$,
$N = 100\,000$ даёт $\hat Q = 0{,}356830$ с ДИ $[0{,}353866;\,
0{,}359805]$ за $1{,}7$~с, и точное значение $0{,}355719$ попадает в
интервал.

\begin{table}[htb]
\centering
\caption{Масштабирование точного перебора (\texttt{connectivity},
$W_{\text{доп}} = 0{,}5$, одинаковые узлы)}
\label{tab:scaling}
\begin{tabular}{@{}rrr@{}}
\toprule
$n$ & $Q(S)$ & Время перебора, с \\
\midrule
10 & 0{,}525681 & 0{,}007 \\
15 & 0{,}449413 & 0{,}279 \\
18 & 0{,}391677 & 2{,}580 \\
20 & 0{,}355719 & 11{,}353 \\
\midrule
\multicolumn{3}{@{}l}{Монте-Карло, $n = 20$, $N = 100\,000$:
$\hat Q = 0{,}356830$,}\\
\multicolumn{3}{@{}l}{ДИ 95\,\% $[0{,}353866;\, 0{,}359805]$, время
$1{,}709$~с}\\
\bottomrule
\end{tabular}
\end{table}

Таким образом, граница $n \leq 20$ для точного перебора обоснована
практическим временем интерактивного сеанса, а при больших $n$ оценка с
доверительным интервалом предпочтительнее и по времени, и по контролируемой
точности~\cite{metropolis,cramer}.

\subsection{Выбор рационального способа защиты}
\label{sec:protection-line}

\textbf{Линейная цепочка (точный перебор).} Для системы
разд.~\ref{sec:protection-line} с бюджетом $C_{\text{доп}} = 50$ рассмотрены
три варианта защиты (табл.~\ref{tab:protection-line}): экранирование ЦОИ
(снижение $p_{ij}$ узла ЦОИ, $C = 25$), резервный канал <<ЦОИ--АРМ>>
(добавление связи, $C = 40$) и комплексная защита (оба мероприятия,
$C = 70$). Живучесть без защиты $Q_{\text{base}}(S) = 0{,}701266$.

\begin{table}[htb]
\centering
\caption{Сравнение способов защиты линейной цепочки ($C_{\text{доп}} = 50$)}
\label{tab:protection-line}
\begin{tabular}{@{}lcclr@{}}
\toprule
Вариант & $C(g)$ & Допустим & $Q_g(S)$ & $\Delta Q$ к базе \\
\midrule
Экранирование ЦОИ & 25 & да & 0{,}712178 & +0{,}010912 \\
\textbf{Резервный канал ЦОИ--АРМ} & 40 & да &
\textbf{0{,}897544} & \textbf{+0{,}196278} \\
Комплексная защита & 70 & нет & 0{,}949679 & +0{,}248412 \\
\bottomrule
\end{tabular}
\end{table}

Рациональный способ~--- резервный канал: $\Delta Q = +0{,}196278$
($+27{,}99\,\%$), эффективность затрат $\Delta Q / C(g^{*}) = 0{,}004907$.
Комплексная защита эффективнее по живучести ($+0{,}248412$), но выходит за
бюджет и отбрасывается до сравнения. Результат согласуется с
разд.~3.3: резервная связь устраняет именно ту потерю состояния $(101)$,
которая ограничивала живучесть линии.

\textbf{Кольцевая сеть из 22 сегментов (Монте-Карло).} Рассмотрим
кольцевую топологию из 22 одинаковых сегментов ($\lambda_i^{(0)} =
0{,}02$, $\mu = 2$, $\Lambda_i = 0{,}245$, $Q_i = 0{,}902429$,
\texttt{connectivity}, $W_{\text{доп}} = 0{,}5$, $N = 50\,000$,
\texttt{seed} $= 42$). Здесь $n > 20$, все оценки~--- Монте-Карло на
общем потоке случайных чисел (табл.~\ref{tab:protection-ring}).

\begin{table}[htb]
\centering
\caption{Сравнение способов защиты кольцевой сети ($C_{\text{доп}} = 50$,
$Q_{\text{base}}(S) = 0{,}517920$)}
\label{tab:protection-ring}
{\footnotesize
\begin{tabular}{@{}lclllr@{}}
\toprule
Вариант & $C(g)$ & Допустим & $\hat Q_g(S)$ & ДИ 95\,\% & $\Delta Q$ \\
\midrule
СОВ и фильтрация & 60 & нет & 0{,}929480 &
[0{,}927203; 0{,}931691] & +0{,}411560 \\
\textbf{Резервные хорды (4 канала)} & 30 & да & \textbf{0{,}877840} &
[0{,}874941; 0{,}880681] & \textbf{+0{,}359920} \\
Ускоренное восстановление & 45 & да & 0{,}869760 &
[0{,}866781; 0{,}872682] & +0{,}351840 \\
\bottomrule
\end{tabular}}
\end{table}

Рациональный способ~--- резервные хорды (четыре дополнительных канала,
$C = 30$): $\Delta Q = +0{,}359920$ ($+69{,}49\,\%$ к базе),
эффективность затрат $\Delta Q / C(g^{*}) = 0{,}011997$. Вариант
<<СОВ и фильтрация>> даёт больший прирост ($+0{,}411560$), но при $C = 60$
недопустим по бюджету; вариант <<ускоренное восстановление>>
($\mu = 6$ на всех сегментах, $C = 45$) уступает хордам и по $\Delta Q$,
и по эффективности затрат ($0{,}351840/45 = 0{,}007819$). Доверительные
интервалы вариантов не перекрываются у выбранного и уступающего вариантов
($[0{,}874941; 0{,}880681]$ против $[0{,}866781; 0{,}872682]$), а оценки
получены на общем потоке чисел, поэтому их различие статистически
устойчиво и воспроизводимо при \texttt{seed} $= 42$.

Заметим, что в обоих примерах рациональным оказывается \emph{структурное}
мероприятие (добавление связей), а не снижение интенсивностей поражения:
топологическая избыточность для малых сетей выигрывает у параметрической
защиты при том же бюджете.

\section{Обсуждение и заключение}

Предложенная методика объединяет три уровня анализа в один воспроизводимый
поток: марковская живучесть узла с совместным учётом отказов и атак
через прореживание пуассоновских потоков; живучесть системы по топологии
с функцией эффективности и критичностями Бирнбаума; выбор рациональной
защиты максимизацией живучести производных систем при бюджетном
ограничении на общем потоке случайных чисел. Численные примеры подтверждают
работоспособность и практическую ценность каждого уровня: доли $s_j$
ранжируют угрозы (71\,\% ущерба узла~--- DDoS); точный перебор на $n = 3$
совпал с ручным расчётом до шестого знака; одна резервная связь подняла
$Q(S)$ линейной цепочки на 28\,\%; для кольцевой сети из 22 сегментов
Монте-Карло с интервалом Уилсона за $1{,}7$~с дал оценку, сопоставимую с
точным перебором; сравнение вариантов защиты выявило рациональный вариант
с эффективностью затрат $0{,}012$ на единицу стоимости.

Ограничения методики определяются её допущениями. Во-первых, состояния
узлов независимы и потоки воздействий пуассоновские с постоянными
интенсивностями: коррелированные атаки (одна кампания бьёт по нескольким
узлам) методика не моделирует напрямую; возможна консервативная оценка
суммарной интенсивностью общих воздействий. Во-вторых, восстановление
экспоненциально с интенсивностью $\mu$; немарковское восстановление
(детерминированное время ремонта, логистика запчастей) аппроксимируется
$\mu = 1 / \bar t_{\text{восст}}$, а для точного учёта требует
имитационного моделирования~\cite{ross}. В-третьих, связи $L$
предполагаются надёжными: отказы каналов моделируются введением каналов
как отдельных узлов. В-четвёртых, точный перебор ограничен $n \leq 20$,
а показатели важности при Монте-Карло оцениваются грубее самой $Q(S)$
($2\,000$ условных прогонов на узел). Наконец, как и всякая аналитическая
модель, методика в первую очередь точна в \emph{сравнительной} картине
(какой узел важнее, какой вариант защиты эффективнее); абсолютные
значения $Q_i$ и $Q(S)$ корректны при калибровке $\lambda_j$ и $p_{ij}$
по статистике конкретного контура.

Направления развития: кусочно-постоянные интенсивности $\lambda_j(t)$
для атак с эскалацией; учёт коррелированных воздействий совместными
марковскими моделями; параллельный Монте-Карло с фиксированным
распределением зёрен; имитационные расширения для немарковского
восстановления; расширение библиотеки ролей узлов графической оболочки.

Тем не менее уже в текущем виде инструмент закрывает практический цикл
<<модель~--- расчёт~--- обоснование решения>>: любая цифра воспроизводится
одной командой по файлу проекта, что делает результаты пригодными для
аудита, отчётности и согласования бюджетов защиты.

\begin{thebibliography}{99}

\bibitem{gnedenko}
Гнеденко Б.~В. Математические методы в теории надёжности /
Б.~В.~Гнеденко, Ю.~К.~Беляев, А.~Д.~Соловьев. ---
М.: Наука, 1965. --- 524~с.

\bibitem{barlow}
Барлоу Р. Математическая теория надёжности /
Р.~Барлоу, Ф.~Прошан ; пер. с англ. ---
М.: Наука, 1969. --- 488~с.

\bibitem{ventcel}
Вентцель Е.~С. Теория вероятностей и её инженерные приложения /
Е.~С.~Вентцель, Л.~А.~Овчаров. ---
М.: Наука, 1988. --- 480~с.

\bibitem{cramer}
Крамер Г. Математические методы статистики /
Г.~Крамер ; пер. с англ. ---
М.: Мир, 1975. --- 648~с.

\bibitem{gost}
ГОСТ 27.002-2015. Надёжность в технике. Термины и определения. ---
М.: Стандартинформ, 2016.

\bibitem{birnbaum}
Birnbaum Z.~W. On the importance of different components in a
multicomponent system / Z.~W.~Birnbaum //
Multivariate Analysis~---~II / ed. P.~R.~Krishnaiah. ---
New York: Academic Press, 1969. --- P.~581--592.

\bibitem{wilson}
Wilson E.~B. Probable inference, the law of succession, and statistical
inference / E.~B.~Wilson // Journal of the American Statistical
Association. --- 1927. --- Vol.~22, №~158. --- P.~209--212.

\bibitem{metropolis}
Metropolis N. The Monte Carlo method / N.~Metropolis, S.~Ulam //
Journal of the American Statistical Association. --- 1949. ---
Vol.~44, №~247. --- P.~335--341.

\bibitem{glasserman}
Glasserman P. Monte Carlo Methods in Financial Engineering /
P.~Glasserman. --- New York: Springer, 2003. --- 596~p. ---
(Stochastic Modelling and Applied Probability; Vol.~53).

\bibitem{rausand}
Rausand M. System Reliability Theory: Models, Statistical Methods, and
Applications / M.~Rausand, A.~H{\o}yland. --- 2nd~ed. ---
Hoboken: Wiley, 2004. --- 644~p.

\bibitem{trivedi}
Trivedi K.~S. Probability and Statistics with Reliability, Queuing, and
Computer Science Applications / K.~S.~Trivedi. --- 2nd~ed. ---
New York: Wiley, 2002. --- 830~p.

\bibitem{ross}
Ross S.~M. Simulation / S.~M.~Ross. --- 5th~ed. ---
Amsterdam: Academic Press, 2012. --- 310~p.

\end{thebibliography}

\end{document}
